\documentclass[11pt]{article}

% ─── Packages ───────────────────────────────────────────────────────────────
\usepackage[margin=1in]{geometry}
\usepackage{amsmath,amssymb,amsthm}
\usepackage{graphicx}
\usepackage{booktabs}
\usepackage{hyperref}
\usepackage{natbib}
\usepackage{xcolor}
\usepackage{bm}

\newtheorem{theorem}{Theorem}
\newtheorem{proposition}{Proposition}
\newtheorem{definition}{Definition}

\hypersetup{
    colorlinks=true,
    linkcolor=blue,
    citecolor=blue,
    urlcolor=blue,
}

% ─── Macros ─────────────────────────────────────────────────────────────────
\newcommand{\vech}{\bm{h}}
\newcommand{\anchor}[1]{A_{#1}}
\newcommand{\barrier}{\tau}
\newcommand{\dive}{D}
\newcommand{\nhat}{\hat{\bm{n}}}
\newcommand{\R}{\mathbb{R}}
\newcommand{\lse}{\mathrm{logsumexp}}

\title{A Trained Iterative Classifier Reduces to Energy Gradient Descent:\\Recovering Closed-Form Dynamics from a Learned NLI System}

\author{
    Chetan Patil \\
    \texttt{chetan12patil@gmail.com} \\
    \url{https://github.com/chetanxpatil/livnium}
}

\date{March 2026}

\begin{document}
\maketitle

% ═══════════════════════════════════════════════════════════════════════════
\begin{abstract}
We train an iterative classification head for natural language inference
in which a hidden state evolves through multiple update steps before
label readout. The update is implemented as a learned MLP (1.2M parameters),
trained end-to-end on SNLI.

We find that the learned update is unnecessary. Replacing the trained MLP
with the analytical gradient of an energy function over anchor cosine
similarities,
$V(\vech) = -\lse(\beta \cdot \cos(\vech, \anchor{k}))$,
produces identical performance on the full SNLI dev set
(82.21\% vs 82.05\%; $n=9{,}842$) without retraining. The trained
network had converged to performing gradient descent on a simple,
interpretable energy landscape---a result verified by ablation, not
assumed by design.

We formalize this as three empirical patterns: (1) the initial state
encodes the relational difference $h_0 = v_q - v_p$ (a design choice),
(2) the representation space behaves as a logsumexp energy over anchor
similarities (discovered), and (3) inference corresponds to gradient
descent on that energy (discovered). Joint retraining under the
recovered dynamics improves neutral recall by $+6.1$pp across five
epochs and reduces head--dynamics disagreement from 55\% to 4\%.
Trajectory analysis reveals a universal fixed-point collapse after
${\sim}3$ steps, identifying the primary accuracy bottleneck.
Code and pretrained weights are publicly available.
\end{abstract}

% ═══════════════════════════════════════════════════════════════════════════
\section{Introduction}
\label{sec:intro}

Standard neural classifiers map a representation $h$ to logits in a
single step. In this work, we consider an alternative: evolving $h$
through a sequence of learned updates before classification, guided by
anchor vectors that define label-specific basins in representation space.

The key finding is that the learned update is not required.

We train such a system on SNLI with a 1.2M-parameter MLP as the update
rule. After training, we replace the MLP with the analytical gradient of
an energy function defined over cosine similarities to the label anchors:
\[
    V(h) = -\log\!\sum_{k} \exp\!\big(\beta \cdot \cos(h, \anchor{k})\big),
    \qquad h_{t+1} = h_t - \alpha \nabla_h V(h_t)
\]
This parameter-free substitution produces the same accuracy on the full
SNLI dev set (82.21\% vs 82.05\%, $n=9{,}842$) without any retraining.
The trained system had converged to performing gradient descent on a
simple, interpretable energy landscape.

This was not designed. It was discovered by systematic ablation: removing
the learned component and checking whether a closed-form alternative
could reproduce the behavior. It could.

\paragraph{Context.} The original Livnium system \citep{patil2026livnium}
introduced attractor-based classification with a bag-of-words encoder
(76\% on SNLI), using hand-designed cosine-radial forces with a
135$^{\circ}$ mismatch from the true cosine gradient---a non-conservative
force field with no scalar potential. Upgrading to a jointly-trained BERT
bi-encoder (82\%) and analyzing the resulting dynamics revealed that the
trained MLP had learned to correct this geometric inconsistency, producing
behavior equivalent to gradient descent on $V(h)$.

\paragraph{Three empirical patterns.} We formalize the discovered behavior as:
\begin{enumerate}
    \item \textbf{Relational initialization} (design choice): $h_0 = v_q - v_p$
    \item \textbf{Energy landscape} (discovered): $V(h) = -\lse\big(\beta \cdot \cos(h, \anchor{E}),\; \beta \cdot \cos(h, \anchor{C}),\; \beta \cdot \cos(h, \anchor{N})\big)$
    \item \textbf{Gradient dynamics} (discovered): $h_{t+1} = h_t - \alpha \nabla V(h_t)$
\end{enumerate}

Patterns 2 and 3 were recovered empirically from a trained system, not
assumed by design. Pattern 1 is a design choice that works well but is
not fundamental---any relational encoding of the inputs should suffice.

\paragraph{Contributions.}
\begin{enumerate}
    \item Demonstration that a trained iterative classifier can be reduced
          to a parameter-free energy gradient with no accuracy loss
          (Section~\ref{sec:discovery}).
    \item Characterization of the recovered dynamics as three empirical
          patterns (Section~\ref{sec:laws}).
    \item Joint retraining framework that aligns the classification head
          with the discovered dynamics (Section~\ref{sec:joint}).
    \item Tunnel test diagnostic that classifies error modes by
          trajectory analysis (Section~\ref{sec:tunnel}).
    \item Identification of a universal fixed-point collapse as the
          primary accuracy bottleneck (Section~\ref{sec:fixedpoint}).
\end{enumerate}

% ═══════════════════════════════════════════════════════════════════════════
\section{Background and Related Work}
\label{sec:related}

\textbf{Energy-based models.}
The idea of classification via energy minimization has a long history.
Hopfield networks \citep{hopfield1982} used energy-based attractor
dynamics for content-addressable memory; modern Hopfield networks
\citep{ramsauer2020hopfield} extend this with exponential storage
capacity and connections to attention mechanisms. LeCun et al.'s
energy-based learning framework \citep{lecun2006tutorial} provides a
general formulation. Our work differs in that the energy function was
not designed but \emph{discovered} from trained dynamics.

\textbf{Equilibrium and implicit models.}
Deep equilibrium models \citep{bai2019deep} compute representations as
fixed points of implicit layers, and neural ODEs \citep{chen2018neural}
parameterize continuous-time dynamics. Livnium uses explicit, finite-step
dynamics with an interpretable geometric structure---the energy landscape
can be written in closed form.

\textbf{Prototype-based classification.}
Prototypical networks \citep{snell2017prototypical} classify by distance
to class centroids. Livnium's anchors serve a similar role, but the
classification involves iterative state refinement, not a single distance
comparison. The anchors define an energy landscape, not just a nearest-neighbor rule.

\textbf{Natural language inference.}
SNLI \citep{bowman2015snli} is the standard benchmark, with baselines
from bag-of-words (${\sim}80\%$) through decomposable attention
\citep{parikh2016decomposable} (${\sim}86\%$) to fine-tuned BERT
\citep{devlin2019bert} (${\sim}90\%$). Our work focuses on the
classification head rather than the encoder; we demonstrate that the
head itself contains discoverable physical laws.

% ═══════════════════════════════════════════════════════════════════════════
\section{The Livnium System}
\label{sec:system}

\subsection{Architecture Overview}

Given a premise $p$ and hypothesis $q$, a BERT bi-encoder produces
sentence vectors $\bm{v}_p, \bm{v}_q \in \R^{768}$. The initial state
for the collapse dynamics is their difference (Law~1):
\[
    h_0 = \bm{v}_q - \bm{v}_p
\]

Three learned anchor vectors $\anchor{E}, \anchor{C}, \anchor{N} \in \R^{768}$
(entailment, contradiction, neutral) define the geometry of the label
space. Each anchor is unit-normalized during the forward pass.

\subsection{Trained Collapse Rule (v1)}

The original system evolves the state through $L=6$ steps using:
\begin{equation}
\label{eq:trained_update}
    h_{t+1} = h_t
    + \delta_\theta(h_t)
    - s_y \cdot D(h_t, \anchor{y}) \cdot \nhat(h_t, \anchor{y})
    - \beta \cdot B(h_t) \cdot \nhat(h_t, \anchor{N})
\end{equation}
where $\delta_\theta$ is a learned 2-layer MLP, $D(h, A) = 0.38 - \cos(h, A)$
is the divergence from the equilibrium ring, $\nhat(h, A) = (h-A)/\|h-A\|$
is the Euclidean radial direction, and $B(h) = 1 - |\cos(h, \anchor{E}) - \cos(h, \anchor{C})|$
is the boundary proximity. The equilibrium target for each anchor is a
cosine-similarity ring at $\cos(h, A_y) = 0.38$, not the anchor point
itself.

\subsection{Classification Head}

After $L$ collapse steps, the refined state $h_L$ is passed to a readout
MLP (SNLIHead) along with auxiliary features: the original premise and
hypothesis vectors $\bm{v}_p, \bm{v}_q$, cosine alignments, and
geometric signals. The head outputs logits for three SNLI classes.

% ═══════════════════════════════════════════════════════════════════════════
\section{The Three Laws}
\label{sec:laws}

We propose that the trained system's behavior is fully specified by
three laws. Law~1 is a design choice; Laws~2 and~3 are empirical
discoveries.

\begin{definition}[Law 1: Relational State Formation]
\begin{equation}
    h_0 = \bm{v}_q - \bm{v}_p
\end{equation}
The initial state encodes the \emph{direction of semantic change} from
premise to hypothesis. Entailment pulls $h_0$ toward $\anchor{E}$;
contradiction toward $\anchor{C}$; neutral is orthogonal.
\end{definition}

\begin{definition}[Law 2: Energy Landscape]
\begin{equation}
\label{eq:energy}
    V(h) = -\log\!\sum_{k \in \{E,C,N\}} \exp\!\big(\beta \cdot \cos(h, \anchor{k})\big)
\end{equation}
The energy at any point $h$ is the negative log-partition function of
Boltzmann-weighted cosine similarities to the three class anchors.
$V(h)$ assigns lower energy to points near an anchor, defining three
semantic basins. The parameter $\beta$ controls basin sharpness.
\end{definition}

\begin{definition}[Law 3: Collapse Dynamics]
\begin{equation}
\label{eq:collapse}
    h_{t+1} = h_t - \alpha \nabla_h V(h_t)
\end{equation}
The state evolves by gradient descent on the energy landscape over $L$
steps. The gradient is computed analytically:
\begin{equation}
\label{eq:gradient}
    \nabla_h V(h) = -\sum_{k} w_k(\vech) \cdot \nabla_h \cos(h, \anchor{k})
\end{equation}
where $w_k(h) = \text{softmax}\big(\beta \cdot \cos(h, \anchor{k})\big)$
are the Boltzmann weights. The cosine gradient for each anchor is:
\begin{equation}
    \nabla_h \cos(h, A) = \frac{A - h \cdot \cos(h, A)}{\|h\|}
\end{equation}
\end{definition}

\paragraph{Key property.} The Boltzmann weights $w_k$ ensure smooth
interpolation near basin boundaries. At $\beta = 1$, the weights form
a balanced softmax over anchors---preserving multi-basin gradient signal
at boundaries. At $\beta \to \infty$, the dynamics reduce to argmax
(nearest-anchor), losing boundary sensitivity.

% ═══════════════════════════════════════════════════════════════════════════
\section{Discovering the Equation of Motion}
\label{sec:discovery}

The three laws were not designed. They were discovered by a systematic
comparison of the trained system with analytical alternatives.

\subsection{Experimental Setup}

We train a BERT bi-encoder jointly with the collapse engine on SNLI
(549k training pairs, 9,842 dev pairs), achieving 82.06\% dev accuracy.
We then evaluate three collapse modes on the full dev set:

\begin{enumerate}
    \item \textbf{Full}: The trained system (Eq.~\ref{eq:trained_update})---MLP $\delta_\theta$ + anchor forces.
    \item \textbf{No-delta}: Remove the learned MLP; keep only anchor forces.
    \item \textbf{Grad-V}: Replace the entire collapse rule with Eq.~\ref{eq:collapse}---pure gradient descent on $V(h)$.
\end{enumerate}

\subsection{Result: The MLP Is Redundant}

\begin{table}[ht]
\centering
\begin{tabular}{@{}lcccc@{}}
\toprule
Mode & Accuracy & E-recall & N-recall & C-recall \\
\midrule
\multicolumn{5}{c}{\textit{Original checkpoint ($\beta=1$, full SNLI dev, $n=9{,}842$)}} \\
\midrule
Full ($\delta_\theta$ + forces) & 82.05\% & 92.79\% & 71.16\% & 81.88\% \\
No-delta (forces only) & 81.81\% & 93.48\% & 68.62\% & 82.98\% \\
\textbf{Grad-V ($\beta=1$, $\alpha=0.2$)} & \textbf{82.21\%} & 92.97\% & \textbf{72.18\%} & 81.18\% \\
\midrule
\multicolumn{5}{c}{\textit{Joint-retrained checkpoint ($\beta=20$, $n=2{,}000$)}} \\
\midrule
Full ($\delta_\theta$ + forces) & 82.35\% & 90.35\% & 71.64\% & 85.30\% \\
No-delta (forces only) & 82.20\% & 93.36\% & 69.87\% & 83.64\% \\
\textbf{Grad-V ($\beta=20$, $\alpha=0.2$)} & \textbf{82.10\%} & \textbf{94.42\%} & \textbf{72.53\%} & 79.55\% \\
\bottomrule
\end{tabular}
\caption{Gradient collapse comparison. Grad-V matches or exceeds the
trained system on overall accuracy and neutral recall. The learned MLP
$\delta_\theta$ contributes $\leq 0.25\%$---it is functionally redundant.
The system \emph{is} a gradient flow on $V(h)$.}
\label{tab:collapse}
\end{table}

The key finding is in Table~\ref{tab:collapse}: the analytical gradient
of Eq.~\ref{eq:energy} matches the trained system. On the full dev set
($n=9{,}842$), Grad-V actually \emph{outperforms} the trained system by
+0.16pp overall and +1.02pp on neutral recall. The learned MLP
$\delta_\theta$, which constitutes 1.2M parameters and requires a
forward pass at each collapse step, contributes at most 0.24pp. It is
functionally dead.

\paragraph{Interpretation.}
The trained MLP was approximating $\nabla V$ but introducing small
distortions. The clean analytical gradient recovers and slightly exceeds
the learned behavior. The energy $V(h)$ is the Boltzmann log-partition
function over anchor similarities---a natural measure of semantic basin
proximity. The system learned to perform gradient descent on this energy
without being told to.

\subsection{Stability Across $\beta$}

We swept $\beta \in \{0.5, 1, 2, 5, 10, 20, 50\}$ at $\alpha = 0.2$.
All values within $\beta \in [1, 50]$ produce accuracy within $\pm 1\%$
of the trained system, confirming that the energy landscape formulation
is robust to the sharpness parameter. The result is not an artifact of
hyperparameter tuning.

% ═══════════════════════════════════════════════════════════════════════════
\section{Joint Retraining Under Discovered Dynamics}
\label{sec:joint}

The Grad-V discovery raises a natural question: if the dynamics are
gradient descent on $V(h)$, can we retrain the system \emph{explicitly}
under these dynamics so that the classification head learns to trust the
energy-based collapse?

\subsection{Motivation: The TYPE-? Problem}

At $\beta=20$ (sharp basins), the gradient dynamics become highly
decisive---states fall cleanly into the nearest anchor basin. However,
the classification head (SNLIHead) was trained under the original
$\beta=1$ dynamics, where collapse produces a universal fixed point
(Section~\ref{sec:fixedpoint}). When we swap in $\beta=20$ grad-V
dynamics at test time, the head's predictions disagree with the
dominant basin on 55.1\% of errors---the head ignores the dynamics it
was not trained on.

\subsection{Method}

We introduce \textbf{LiveniumJoint}, a joint training framework:
\begin{itemize}
    \item BERT bi-encoder, standalone anchors $[\anchor{E}, \anchor{C}, \anchor{N}] \in \R^{3 \times 768}$, and SNLIHead are trained together.
    \item Collapse uses differentiable grad-V dynamics (Eq.~\ref{eq:collapse}) with analytical gradient, staying in the computation graph.
    \item Loss combines cross-entropy with an anchor alignment term:
\end{itemize}
\begin{equation}
\label{eq:joint_loss}
    \mathcal{L} = \mathcal{L}_{\text{CE}} + \lambda_{\text{align}} \cdot \big(-\cos(h_L, \anchor{y})\big)
\end{equation}

The alignment loss directly penalizes disagreement between the final
collapsed state and the correct anchor, forcing the head to be
consistent with the dynamics.

\subsection{Results}

We train on 30,000 stratified samples (10k per class) for 5 epochs
with $\beta=20$, $\alpha=0.05$, $\lambda_{\text{align}}=0.3$, warm-starting
from the BERT-joint checkpoint.

\begin{table}[ht]
\centering
\begin{tabular}{@{}lcccccc@{}}
\toprule
Metric & Before & Ep.\ 1$^\star$ & Ep.\ 2 & Ep.\ 3 & Ep.\ 4 & Ep.\ 5 \\
\midrule
Accuracy       & 81.95\% & \textbf{82.79\%} & 82.48\% & 82.68\% & 82.43\% & 82.48\% \\
Neutral recall & 70.54\% & 74.50\% & 74.74\% & 75.98\% & \textbf{76.63\%} & 76.51\% \\
C recall       & ---     & 80.96\% & 79.62\% & 80.32\% & 78.83\% & 79.23\% \\
Dyn/head agree & 29.73\% & 72.52\% & 74.10\% & 73.29\% & 73.62\% & 73.88\% \\
TYPE-? rate    & 7.84\%  & 4.20\%  & 3.98\%  & 4.56\%  & 4.50\%  & 4.43\%  \\
\midrule
$\cos(\anchor{C},\anchor{N})$ & --- & $-0.031$ & $-0.007$ & $+0.011$ & $+0.024$ & $+0.026$ \\
\bottomrule
\end{tabular}
\caption{Full 5-epoch joint retraining ($^\star$ = best checkpoint).
Neutral recall improves \emph{monotonically} across all epochs
($70.5\% \to 76.6\%$, $+6.1$pp), peaking at epoch~4.
Overall accuracy is best at epoch~1 (82.79\%), reflecting a trade-off
with C-recall. The final row shows $\cos(\anchor{C},\anchor{N})$
crossing zero at epoch~3 and reaching $+0.026$ by epoch~5---confirmed
C--N anchor drift.}
\label{tab:joint}
\end{table}

Joint retraining achieves its primary goal. Head--dynamics agreement
rises from 29.7\% to 73.9\%, and neutral recall---the weakest
class---improves monotonically across all five epochs, from 70.54\%
to 76.63\% ($+6.1$pp). Overall accuracy is best at epoch~1 (82.79\%),
reflecting a trade-off: the alignment loss optimizes neutral recall at
some cost to contradiction recall and overall accuracy.
The Grad-V evaluation on the best checkpoint
(Table~\ref{tab:collapse}, bottom) confirms that the analytical gradient
still matches the full system, with N-recall of 72.53\% exceeding the
full model's 71.64\%.

\paragraph{C--N anchor drift.}
$\cos(\anchor{C}, \anchor{N})$ crosses zero at epoch~3 and reaches
$+0.026$ by epoch~5. The C and N anchors become \emph{positively}
correlated under the alignment loss, and C-recall erodes from 80.96\%
to 78.83\% correspondingly. This identifies a concrete regularization
target: anchor orthogonality must be explicitly enforced during joint
training. Adding a penalty $\lambda_{\text{orth}}\sum_{i \neq j}\cos(\anchor{i},\anchor{j})^2$
is the natural fix and a direct target for future work.

% ═══════════════════════════════════════════════════════════════════════════
\section{Trajectory Analysis}
\label{sec:analysis}

\subsection{Tunnel Test: Classifying Error Modes}
\label{sec:tunnel}

We develop a \textbf{tunnel test} diagnostic that traces the step-by-step
trajectory of each collapse and classifies errors by \emph{when} and
\emph{how} the dynamics fail. For each sample, we record the dominant
anchor basin (highest cosine) at every step $t = 0, \ldots, L$.

Error types for neutral misclassifications:
\begin{itemize}
    \item \textbf{TYPE-1} (bad $h_0$, Law~1 failure): $h_0$ starts in the wrong basin. The initial encoding does not place the state near the neutral anchor.
    \item \textbf{TYPE-2} (mid-diversion, Law~3 failure): $h_0$ starts correctly but the trajectory is diverted to a wrong basin during collapse.
    \item \textbf{TYPE-3} (boundary stall, Law~2 failure): the trajectory converges near a basin boundary or to the universal fixed point, never reaching a decisive basin.
    \item \textbf{TYPE-?} (head override): the dynamics deliver the correct dominant basin at step $L$, but the classification head overrides and predicts a different label.
\end{itemize}

\begin{table}[ht]
\centering
\begin{tabular}{@{}lrrrr@{}}
\toprule
& TYPE-1 & TYPE-2 & TYPE-3 & TYPE-? \\
\midrule
$\beta=1$ (original) & 14.4\% & 21.3\% & \textbf{63.9\%} & 0.5\% \\
$\beta=20$ (grad-V) & 2.8\% & 20.4\% & 21.7\% & \textbf{55.1\%} \\
$\beta=20$ (joint) & --- & --- & --- & \textbf{4.0\%} \\
\bottomrule
\end{tabular}
\caption{Error type distribution across configurations. At $\beta=1$,
TYPE-3 (fixed-point convergence) dominates. At $\beta=20$, dynamics
become decisive but the head overrides them (TYPE-?). Joint retraining
resolves the head--dynamics disconnect.}
\label{tab:tunnel}
\end{table}

\subsection{The Universal Fixed Point}
\label{sec:fixedpoint}

The dominant error mode at $\beta=1$ is TYPE-3 (63.9\%). Investigation
reveals this is not boundary ambiguity but a more fundamental phenomenon:
the collapse dynamics converge to a \textbf{universal fixed point} after
approximately 3 steps, independent of the input.

At steps 4--6, nearly all samples---correct and incorrect---produce
the same cosine alignment pattern:
\[
    \cos(h, \anchor{E}) \approx -0.129, \quad
    \cos(h, \anchor{C}) \approx -0.124, \quad
    \cos(h, \anchor{N}) \approx +0.172
\]

This fixed point is weakly neutral-biased ($\cos(h, \anchor{N})$ is
highest), explaining why neutral recall is better than chance (71\%)
but not as strong as entailment or contradiction. The collapse discards
input-specific information after step 3---the final state is
input-agnostic.

\paragraph{Energy signature.}
Correctly classified samples show mean energy descent
$\Delta V = +0.046$ (ascending---moving toward lower-energy basins),
while error samples show $\Delta V = -0.002$ (flat---stuck at the
fixed point). The dynamics are contracting for correct predictions and
stagnant for errors.

\paragraph{Implication.}
The accuracy bottleneck is not the classification head or the energy
landscape---it is the fixed-point collapse. Input information is lost
after three steps. Breaking this fixed point (e.g., via input-dependent
residual forcing or adaptive step sizes) is the primary target for
future work.

\subsection{Basin Stability}
\label{sec:stability}

Correct predictions occupy deeper attractor basins than incorrect ones.
We measure stability by perturbing $h_0$ with Gaussian noise
($\sigma = 0.3$) across 20 trials per sample:

\begin{table}[ht]
\centering
\begin{tabular}{@{}lcc@{}}
\toprule
Group & Flip rate & Entropy \\
\midrule
Correct ($n=1{,}653$) & $0.0017 \pm 0.025$ & 0.293 \\
Incorrect ($n=347$) & $0.0068 \pm 0.048$ & 0.604 \\
\textbf{Ratio} & \textbf{3.93$\times$} & \textbf{2.06$\times$} \\
\bottomrule
\end{tabular}
\caption{Basin stability: wrong predictions are 3.9$\times$ more likely
to flip under perturbation and carry 2$\times$ higher entropy.
Correctness correlates with basin depth---no calibration needed.}
\label{tab:stability}
\end{table}

Per-class breakdown: entailment errors flip 7.18$\times$ more than
correct entailment predictions, neutral errors 3.17$\times$,
contradiction errors only 1.82$\times$. The low contradiction ratio
indicates systematic geometric bias rather than boundary ambiguity---a
structural feature of the anchor geometry.

% ═══════════════════════════════════════════════════════════════════════════
\section{Discussion}
\label{sec:discussion}

\subsection{Why Does the MLP Learn $\nabla V$?}

The trained collapse rule (Eq.~\ref{eq:trained_update}) uses
cosine-magnitude forces with Euclidean-radial directions---a
135$^{\circ}$ mismatch from the true cosine gradient. Yet the
\emph{combined} effect of MLP + forces approximates gradient descent
on $V(h)$. We hypothesize that the MLP learns a corrective residual
that compensates for the geometric inconsistency of the hand-designed
forces, effectively rotating the force field toward the true gradient.
The end-to-end training signal (cross-entropy on the final label)
provides sufficient pressure to discover this correction.

\subsection{What Is the Energy $V(h)$?}

The energy $V(h) = -\lse(\beta \cdot \cos(h, \anchor{k}))$ is the
negative log-partition function of a Boltzmann distribution over anchor
cosine similarities. At any point $h$, it measures the aggregate
proximity to all class anchors, weighted exponentially. The minima of
$V$ correspond to anchor positions; the saddle points define basin
boundaries. At $\beta=1$, the Boltzmann weights are a balanced softmax,
preserving gradient signal at boundaries. At $\beta \to \infty$,
$V(h) \to -\max_k \cos(h, \anchor{k})$, recovering the Voronoi
partition over anchors.

\subsection{Accuracy Context}

The system achieves 82--83\% on SNLI, below the BERT-linear baseline
(${\sim}90\%$). The gap is attributable to two factors: (1) the
bi-encoder architecture (premise and hypothesis encoded separately)
loses cross-attention signal that a BERT cross-encoder exploits, and
(2) the universal fixed-point collapse discards input information after
3 steps. The contribution of this work is not state-of-the-art accuracy
but the \emph{discovery of physical laws} governing a neural inference
system---a result that holds regardless of absolute accuracy level.

\subsection{Limitations}

\begin{enumerate}
    \item \textbf{Single dataset.} All results are on SNLI. Evaluation
          on MultiNLI, adversarial NLI, and non-NLI classification tasks
          would test generality.
    \item \textbf{Single seed.} The gradient-flow result has been
          validated on the full dev set ($n=9{,}842$) but from a single
          training run. Multi-seed validation is needed.
    \item \textbf{Encoder dependence.} The laws are demonstrated with a
          BERT bi-encoder. Whether they hold for cross-encoders or
          non-BERT architectures is untested.
    \item \textbf{Fixed-point bottleneck.} The universal fixed point
          limits accuracy. Breaking it requires architectural changes
          not yet validated.
    \item \textbf{C--N anchor drift.} Joint retraining shows signs of
          C--N anchor convergence, which may limit contradiction recall.
\end{enumerate}

% ═══════════════════════════════════════════════════════════════════════════
\section{Conclusion}
\label{sec:conclusion}

We have shown that a trained attractor-based classification system
for NLI is governed by three empirical laws: relational state formation,
a logsumexp cosine energy landscape, and gradient-descent collapse
dynamics. The central finding is that Laws~2 and~3 were \emph{not
designed}---they were recovered by analyzing what the trained MLP was
doing. The learned residual $\delta_\theta$ can be replaced by the
analytical gradient $\nabla V$ with no accuracy loss on 9,842 samples.

Trajectory analysis via the tunnel test reveals the primary bottleneck:
a universal fixed point that discards input information after three
collapse steps. Joint retraining under the discovered dynamics reduces
head--dynamics disagreement from 55\% to 4\% and improves neutral
recall by +4pp.

The broader implication is that neural systems can develop interpretable
physical laws through training---laws that are simpler, faster, and more
transparent than the trained components they replace. Whether this
phenomenon generalizes beyond attractor-based classification heads is
an open question and a promising direction for future work.

\paragraph{Code and Models.}
\url{https://github.com/chetanxpatil/livnium} \\
\url{https://huggingface.co/chetanxpatil/livnium-snli}

% ═══════════════════════════════════════════════════════════════════════════
\bibliographystyle{plainnat}
\begin{thebibliography}{12}

\bibitem[Bai et~al.(2019)]{bai2019deep}
S.~Bai, J.~Z. Kolter, and V.~Koltun.
\newblock Deep equilibrium models.
\newblock In \emph{NeurIPS}, 2019.

\bibitem[Bowman et~al.(2015)]{bowman2015snli}
S.~Bowman, G.~Angeli, C.~Potts, and C.~Manning.
\newblock A large annotated corpus for learning natural language inference.
\newblock In \emph{EMNLP}, 2015.

\bibitem[Chen et~al.(2018)]{chen2018neural}
R.~T.~Q. Chen, Y.~Rubanova, J.~Bettencourt, and D.~Duvenaud.
\newblock Neural ordinary differential equations.
\newblock In \emph{NeurIPS}, 2018.

\bibitem[Devlin et~al.(2019)]{devlin2019bert}
J.~Devlin, M.-W. Chang, K.~Lee, and K.~Toutanova.
\newblock {BERT}: Pre-training of deep bidirectional transformers for
  language understanding.
\newblock In \emph{NAACL-HLT}, 2019.

\bibitem[Hopfield(1982)]{hopfield1982}
J.~J. Hopfield.
\newblock Neural networks and physical systems with emergent collective
  computational abilities.
\newblock \emph{PNAS}, 79(8):2554--2558, 1982.

\bibitem[LeCun et~al.(2006)]{lecun2006tutorial}
Y.~LeCun, S.~Chopra, R.~Hadsell, M.~Ranzato, and F.~Huang.
\newblock A tutorial on energy-based learning.
\newblock In \emph{Predicting Structured Data}. MIT Press, 2006.

\bibitem[Parikh et~al.(2016)]{parikh2016decomposable}
A.~Parikh, O.~T{\"a}ckstr{\"o}m, D.~Das, and J.~Uszkoreit.
\newblock A decomposable attention model for natural language inference.
\newblock In \emph{EMNLP}, 2016.

\bibitem[Patil(2026)]{patil2026livnium}
C.~Patil.
\newblock Iterative attractor dynamics for classification: A dynamical
  classification head with anchor-guided basin settling.
\newblock Preprint, Zenodo, 2026.
\newblock \url{https://zenodo.org/records/19058910}.

\bibitem[Ramsauer et~al.(2021)]{ramsauer2020hopfield}
H.~Ramsauer, B.~Sch{\"a}fl, J.~Lehner, P.~Seidl, M.~Widrich,
  T.~Adler, L.~Gruber, M.~Holzleitner, M.~Pavlovi{\'c}, G.~Sandve,
  V.~Greiff, D.~Kreil, M.~Kopp, G.~Klambauer, J.~Brandstetter,
  and S.~Hochreiter.
\newblock Hopfield networks is all you need.
\newblock In \emph{ICLR}, 2021.

\bibitem[Snell et~al.(2017)]{snell2017prototypical}
J.~Snell, K.~Swersky, and R.~Zemel.
\newblock Prototypical networks for few-shot learning.
\newblock In \emph{NeurIPS}, 2017.

\end{thebibliography}

\end{document}
